% Options for packages loaded elsewhere
\PassOptionsToPackage{unicode}{hyperref}
\PassOptionsToPackage{hyphens}{url}
\documentclass[
  ignorenonframetext,
]{beamer}
\newif\ifbibliography
\usepackage{pgfpages}
\setbeamertemplate{caption}[numbered]
\setbeamertemplate{caption label separator}{: }
\setbeamercolor{caption name}{fg=normal text.fg}
\beamertemplatenavigationsymbolsempty
% remove section numbering
\setbeamertemplate{part page}{
  \centering
  \begin{beamercolorbox}[sep=16pt,center]{part title}
    \usebeamerfont{part title}\insertpart\par
  \end{beamercolorbox}
}
\setbeamertemplate{section page}{
  \centering
  \begin{beamercolorbox}[sep=12pt,center]{section title}
    \usebeamerfont{section title}\insertsection\par
  \end{beamercolorbox}
}
\setbeamertemplate{subsection page}{
  \centering
  \begin{beamercolorbox}[sep=8pt,center]{subsection title}
    \usebeamerfont{subsection title}\insertsubsection\par
  \end{beamercolorbox}
}
% Prevent slide breaks in the middle of a paragraph
\widowpenalties 1 10000
\raggedbottom
\AtBeginPart{
  \frame{\partpage}
}
\AtBeginSection{
  \ifbibliography
  \else
    \frame{\sectionpage}
  \fi
}
\AtBeginSubsection{
  \frame{\subsectionpage}
}
\usepackage{iftex}
\ifPDFTeX
  \usepackage[T1]{fontenc}
  \usepackage[utf8]{inputenc}
  \usepackage{textcomp} % provide euro and other symbols
\else % if luatex or xetex
  \usepackage{unicode-math} % this also loads fontspec
  \defaultfontfeatures{Scale=MatchLowercase}
  \defaultfontfeatures[\rmfamily]{Ligatures=TeX,Scale=1}
\fi
\usepackage{lmodern}
\ifPDFTeX\else
  % xetex/luatex font selection
\fi
% Use upquote if available, for straight quotes in verbatim environments
\IfFileExists{upquote.sty}{\usepackage{upquote}}{}
\IfFileExists{microtype.sty}{% use microtype if available
  \usepackage[]{microtype}
  \UseMicrotypeSet[protrusion]{basicmath} % disable protrusion for tt fonts
}{}
\makeatletter
\@ifundefined{KOMAClassName}{% if non-KOMA class
  \IfFileExists{parskip.sty}{%
    \usepackage{parskip}
  }{% else
    \setlength{\parindent}{0pt}
    \setlength{\parskip}{6pt plus 2pt minus 1pt}}
}{% if KOMA class
  \KOMAoptions{parskip=half}}
\makeatother
\setlength{\emergencystretch}{3em} % prevent overfull lines
\providecommand{\tightlist}{%
  \setlength{\itemsep}{0pt}\setlength{\parskip}{0pt}}
% Slide layout defaults for warning-clean scientific decks.
\usepackage{etoolbox}
\IfFileExists{xurl.sty}{\usepackage{xurl}}{}
\IfFileExists{seqsplit.sty}{\usepackage{seqsplit}}{\newcommand{\seqsplit}[1]{#1}}
\protected\def\breakseq#1{\seqsplit{#1}}
\protected\def\breaktt#1{\begingroup\ttfamily\seqsplit{#1}\endgroup}
\setlength{\emergencystretch}{6em}
\tolerance=5000
\hbadness=10000
\hfuzz=1pt
\setlength{\tabcolsep}{2pt}
\AtBeginEnvironment{longtable}{\tiny\renewcommand{\arraystretch}{0.86}\setlength{\tabcolsep}{1pt}}
\AtBeginEnvironment{tabular}{\tiny\renewcommand{\arraystretch}{0.86}\setlength{\tabcolsep}{1pt}}

% Natbib and cross-reference fallbacks — slides don't load natbib
% or cleveref, but combined-PDF manuscript prose may emit these
% commands. The fallback renders citations as a bracketed key list
% and cross-references as detokenized labels so slides don't fail on
% undefined control sequences or raw underscores. \providecommand is
% a no-op if packages load later.
\providecommand{\citep}[1]{[#1]}
\providecommand{\citet}[1]{#1}
\providecommand{\citealp}[1]{#1}
\providecommand{\citeauthor}[1]{#1}
\providecommand{\citeyear}[1]{#1}
\providecommand{\cref}[1]{\texttt{\detokenize{#1}}}
\providecommand{\Cref}[1]{\texttt{\detokenize{#1}}}

\newtheorem{proposition}{Proposition}
\newtheorem{hypothesis}{Hypothesis}
\usepackage{bookmark}
\IfFileExists{xurl.sty}{\usepackage{xurl}}{} % add URL line breaks if available
\urlstyle{same}
% fallback for those not using the hyperref driver hyperxmp:
\makeatletter
\@ifundefined{xmpquote}{\newcommand{\xmpquote}[1]{#1}}{}
\makeatother
\hypersetup{
  hidelinks,
  pdfcreator={LaTeX via pandoc}}

\author{\texorpdfstring{}{}}
\date{}

\begin{document}

\section{Conclusion}\label{sec:conclusion}

\begin{frame}[fragile,allowframebreaks]{Conclusion}
\breaktt{template\_search\_project} packages a complete, configurable,
reproducible literature workflow into the Research Project Template's
two-layer architecture. By keeping discovery, export, and synthesis in
three orthogonal infrastructure modules, the project demonstrates that
ambitious research automation can still respect the template's
principles:

\begin{itemize}
\tightlist
\item
  \textbf{Single source of truth} --- \texttt{Paper} for discovery,
  \texttt{BibEntry} for export, structured \texttt{SynthesisResult}
  records for LLM output.
\item
  \textbf{Test-driven development} --- every module is covered by
  real-data tests; HTTP backends are exercised through
  \breaktt{pytest-httpserver}, the LLM bridge through deterministic local
  callables.
\item
  \textbf{Thin orchestrator pattern} ---
  \breaktt{scripts/run\_search\_pipeline.py} does only argument parsing,
  configuration loading, and I/O; all logic lives in
  \texttt{infrastructure/} or \texttt{src/}.
\item
  \textbf{No mocks} --- neither in the new infrastructure modules nor in
  the project test suite.
\item
  \textbf{Multi-project support} --- the project lives alongside
  \breaktt{template\_code\_project/} and follows the same layout, so the
  existing pipeline runner discovers and executes it without
  modification.
\item
  \textbf{Reproducibility} --- deterministic search caching, on-disk
  enrichment caching, and pinned LLM seeds make a single
  \breaktt{manuscript/config.yaml} the only artifact a reviewer needs.
\end{itemize}

We close with three concrete extensions that build naturally on this
foundation:

\begin{enumerate}
\tightlist
\item
  \textbf{Crossref TDM full-text fetch} for non-arXiv DOIs, completing
  the abstract-to-fulltext picture without changing the project's API.
\item
  \textbf{CSL-JSON export} alongside BibTeX, enabling Zotero / Mendeley
  / Pandoc-CSL workflows from the same \texttt{BibDatabase}.
\item
  \textbf{Vector recall on \texttt{LocalBackend}} for curated corpora
  exceeding \textasciitilde1000 papers, gated behind an optional
  dependency.
\end{enumerate}

The infrastructure modules are deliberately small and stable; the
project that exercises them is deliberately small and explicit. Together
they show that \emph{domain-specific research automation} and
\emph{template-strict architectural discipline} are compatible --- and,
in fact, mutually reinforcing.

The bundled \breaktt{data/corpus.json} exercises classical optimisation
references \breakseq{[@boyd2004convex; @nocedal2006numerical;
@nesterov2013gradient]} alongside modern stochastic-optimisation work
\breakseq{[@kingma2014adam; @reddi2018convergence]} and the canonical
reproducibility paper \breakseq{[@peng2011reproducible]}, so the auto-generated
\breaktt{manuscript/references.bib} always contains real citation-ready
entries that downstream tooling can resolve.
\end{frame}

\end{document}
